\section{Observability}
\label{sec:observability}

\vannak{} exposes observability as owned, point-in-time snapshots rather than
borrowed views into live state. This keeps reads cheap, non-blocking, and safe
to move across shard boundaries.

\subsection{Snapshot Types}

\subsubsection{HotIndexSnapshot}

\texttt{HotIndexSnapshot} summarizes a shard-local index with counts only: the
number of process instances, events, pipelines, and distinct metadata
references, plus the running \texttt{duplicate\_events} and
\texttt{rejected\_events} counters.

\subsubsection{MetadataOutboxSnapshot}

\texttt{MetadataOutboxSnapshot} reports the outbox queue health.

\begin{lstlisting}
pub struct MetadataOutboxSnapshot {
    pub total: usize,
    pub pending: usize,
    pub acknowledged: usize,
    pub failed: usize,
    pub queued_pending: usize,
}
\end{lstlisting}

The \texttt{failed} count is the primary alerting signal: a sustained nonzero
value means deliveries are not reaching \ipto{}. The gap between \texttt{pending}
and \texttt{queued\_pending} reflects entries that have been dequeued but not yet
resolved.

\subsubsection{SegmentBackedMetadataOutboxSnapshot}

\texttt{SegmentBackedMetadataOutboxSnapshot} pairs the outbox snapshot with the
current \texttt{SegmentManifest}, so a single read shows both the in-memory
delivery state and the durable backing segment.

\begin{lstlisting}
pub struct SegmentBackedMetadataOutboxSnapshot {
    pub outbox: MetadataOutboxSnapshot,
    pub segment: SegmentManifest,
}
\end{lstlisting}

\subsubsection{BoundedIngestRuntimeSnapshot}

\texttt{BoundedIngestRuntimeSnapshot} pairs the aggregate shard runtime snapshot
with one \texttt{IngestQueueSnapshot} per logical shard. Each queue snapshot
reports the \texttt{shard\_id}, current \texttt{depth}, and configured
\texttt{capacity}; sustained depth near capacity is the local backpressure signal
for process-event ingest.

\subsubsection{DataProvenanceIndexSnapshot}

\texttt{DataProvenanceIndexSnapshot} reports recent local provenance lookup
coverage: metadata event count, data-individual count, process-instance count,
and activity count. It is a hot-index observability surface, not the durable
Ipto record of truth.

\subsection{Segment Manifest}

A \texttt{SegmentManifest} is itself an observability surface: it reports the
segment id, owning node, path, \texttt{record\_count}, \texttt{byte\_len}, and
rolling \texttt{checksum}. Comparing record counts and checksums across nodes is
how operators confirm a sealed segment is intact and replicated.

\subsection{Replay Summary}

After a restart, \texttt{MetadataOutboxReplaySummary} reports what replay did:
the \texttt{checkpoint\_offset} used, \texttt{scanned\_records},
\texttt{skipped\_records} (already acknowledged), and \texttt{replayed\_records}
(re-queued as pending). A large skipped count confirms the checkpoint is being
honored.

\subsection{Rebalance Summary}

\texttt{MetadataOutboxRebalanceSummary} reports a shard-range rebalance:
\texttt{entries\_found}, \texttt{attempted}, \texttt{acknowledged}, and
\texttt{failed}. The related \texttt{MetadataOutboxDrainSummary} reports a drain
attempt, including the \texttt{stopped\_after\_failure} flag.

\subsection{Query Types}

The query results from Section~\ref{sec:ingest} double as observability outputs.
\texttt{QueryResult::ProcessInstances} returns owned process snapshots, and
\texttt{QueryResult::Events} returns owned \texttt{PipelineEvent} values for a
process instance or for a metadata-impact lookup.

\subsection{Process Instance Snapshot}

\texttt{ProcessInstanceSnapshot} is the owned per-instance view returned by
queries. It includes the full identity (tenant, environment, pipeline,
definition, instance, version), the current activity and \texttt{ProcessStatus},
start and completion timestamps, correlation and business keys, the per-activity
\texttt{ActivityState} map, per-activity durations, accumulated metadata
references, the \texttt{retry\_count}, and the last error.

\begin{vnote}
Activity durations are computed in milliseconds from the recorded entry time to
the terminal event time. They are derived during reduction, so reading them from
a snapshot involves no recomputation.
\end{vnote}

\subsection{Probe Command}

\texttt{vannak-node probe <host> <port>} opens a connection to a node, sends a
\texttt{ClusterSummaryRequest} envelope, and prints the response payload. It is
the simplest liveness and cluster-view check for a running node.

\begin{lstlisting}[numbers=none]
vannak-node probe localhost 10081
\end{lstlisting}

\subsection{Health Endpoints}

The node exposes its \graft{} transport on port 10081, which serves as the
reachability endpoint probed above. There is no separate HTTP health port; the
transport listener and the \texttt{probe} subcommand together provide the
liveness and summary signals.

\subsection{Container Health Checks}

The test PostgreSQL container defines an explicit health check so dependent
steps can wait for readiness:

\begin{lstlisting}[numbers=none]
healthcheck:
  test: ["CMD-SHELL", "pg_isready -U repo -d repo"]
  interval: 2s
  timeout: 5s
  retries: 30
\end{lstlisting}

For the \raft{} cluster, the integration script checks health by counting
running containers (expecting four: DNS plus three nodes) and probing each node
in turn. A node that fails to respond to \texttt{probe} is reported as
unreachable.

\endinput
